% Build example:
% /mnt/c/cygwin64/bin/pdflatex -interaction=nonstopmode -halt-on-error \
%   -output-directory=/tmp/board_generator_theory_build docs/board_generator_theory.tex
\documentclass[11pt]{article}

\usepackage[a4paper,margin=1in]{geometry}
\usepackage{amsmath,amssymb,bm}

\setlength{\parindent}{0pt}
\setlength{\parskip}{0.6em}

\title{Theory of Knot, Growth-Layer, and Fiber Modeling\\in the Board Generator}
\date{April 17, 2026}

\begin{document}
\maketitle

This report describes the mathematical model used by the current board generator implementation. The
presentation follows the code convention directly:
\[
X = \text{width}, \qquad Y = \text{thickness}, \qquad Z = \text{length}.
\]
The notation is therefore consistent with the updated growth-layer report and with the current
implementation of the simulator. The theory is organized in the same order as the runtime pipeline:
first the knot generator, then the growth-layer generator, and finally the fiber-orientation model.

\section{Knot Generator}

\subsection{Knot Parameters}

Each runtime knot $n$ is described by the parameter vector
\begin{equation}
\kappa_n =
\left(
\theta_{0,n}, z_{0,n}, c_{1,n}, c_{2,n}, k_n, k_{p,n},
A_{\mathrm{bump},n}, B_{\mathrm{bump},n}, A_{\mathrm{exp},n},
a_{1,n}, a_{2,n}, a_{3,n}, a_{4,n}, R_{L,n}, R_{D,n}
\right).
\end{equation}
The roles of these parameters are:
\begin{enumerate}
\item $\theta_{0,n}$ is the azimuth of the knot in the transverse $(X,Y)$ plane.
\item $z_{0,n}$ is the longitudinal anchor position of the knot along the board.
\item $c_{1,n}$ and $c_{2,n}$ are the coefficients of the knot axis parabola.
\item $k_n$ and $k_{p,n}$ control the growth-layer contraction and anisotropy.
\item $A_{\mathrm{bump},n}$, $B_{\mathrm{bump},n}$, and $A_{\mathrm{exp},n}$ define the radial bump law.
\item $a_{1,n},\ldots,a_{4,n}$ are the coefficients of the quartic knot diameter law.
\item $R_{L,n}$ and $R_{D,n}$ are the live and dead transition radii.
\end{enumerate}

The axis and the diameter law associated with one knot are written as
\begin{equation}
Z_{c,n}(\xi) = c_{1,n}\xi^2 + c_{2,n}\xi + z_{0,n},
\end{equation}
\begin{equation}
L_n(\xi) = a_{1,n}\xi^4 + a_{2,n}\xi^3 + a_{3,n}\xi^2 + a_{4,n}\xi.
\end{equation}

The sequence model does not sample the full parameter vector directly. Instead it samples a dictionary row
\begin{equation}
d_m = \left(\theta_{0,m}, R_{L,m}, R_{D,m}, a_{1,m}, a_{2,m}, a_{3,m}, a_{4,m}, c_{1,m}, c_{2,m}\right),
\end{equation}
and the remaining quantities are reconstructed during runtime decoding. In particular, $z_{0,n}$ is obtained
from the longitudinal slot location, while $k_n$, $k_{p,n}$, $A_{\mathrm{bump},n}$, $B_{\mathrm{bump},n}$,
and $A_{\mathrm{exp},n}$ are derived from the sampled diameter profile.

\subsection{Training Dataset}

The knot-sequence training dataset is assembled from raw MATLAB files matching
\texttt{knot\_data\_*.mat}.
For each observed knot row, the preparation stage reads the matrices
\begin{align*}
&\texttt{Knot\_param\_\_Alfa\_L100\_Rlivelim\_Rtot\_Abump\_Aexp},\\
&\texttt{Knot\_param\_\_Vertdiamparam},\\
&\texttt{Knot\_param\_\_Zcparam},
\end{align*}
and constructs the runtime dictionary row $d_m$ from the corresponding columns.

The raw longitudinal start coordinate is read as $z_{0,m}$ and aligned to a regular slot lattice of spacing
$\Delta z$:
\begin{equation}
z^{(\Delta)}_{0,m} = \operatorname{round}\!\left(\frac{z_{0,m}}{\Delta z}\right)\Delta z.
\end{equation}
If two consecutive knots land in the same slot after rounding, the later one is shifted forward by one slot.
Negative slots are discarded.

Each log then becomes a discrete sequence
\begin{equation}
s^{(\ell)} = \left(s^{(\ell)}_1, s^{(\ell)}_2, \ldots, s^{(\ell)}_{T_\ell}\right),
\qquad
s^{(\ell)}_j \in \{0,1,\ldots,M\},
\end{equation}
where $0$ denotes an empty slot and every nonzero symbol refers to one knot dictionary entry. The sequence is
segmented into overlapping windows of length $N_s = 101$ with overlap $97$, so the stride is four slots. For
each window
\begin{equation}
w = \left(t_1, t_2, \ldots, t_{101}\right),
\end{equation}
the supervised training pair is
\begin{equation}
\text{input} = \left(t_1,\ldots,t_{100}\right), \qquad
\text{target} = \left(t_2,\ldots,t_{101}\right).
\end{equation}

In addition to the runtime dictionary, the data-preparation code stores an auxiliary 23-dimensional shape
descriptor. The first three entries are $(\theta_0,R_L,R_D)$; the next ten entries are samples of the diameter
law along normalized axial positions; and the final ten entries are samples of the axis law along normalized
positions. This auxiliary representation is retained for analysis and compatibility. The actual sequence model
still operates on the tokenized knot sequences.

\subsection{Clustering of Knots in the Training Set}

Two tokenization modes are supported.

In the direct mode, each nonzero token corresponds to exactly one dictionary row. In the clustered mode, the
dictionary rows are first normalized columnwise to $[0,1]$:
\begin{equation}
\widehat{d}_m = \frac{d_m - d_{\min}}{d_{\max} - d_{\min}},
\end{equation}
and the nonzero rows are grouped into $Q$ clusters by $k$-means. The clustering problem is
\begin{equation}
\min_{\{\mu_q\}, \{c_m\}}
\sum_{m=1}^{M}
\left\lVert \widehat{d}_m - \mu_{c_m} \right\rVert_2^2,
\qquad
c_m \in \{1,\ldots,Q\},
\end{equation}
where $\mu_q$ is the centroid of cluster $q$ and $c_m$ is the assigned cluster index of dictionary row $m$.

Token $0$ is always reserved for the no-knot state. Therefore the learned vocabulary is
\begin{equation}
\mathcal{V} = \{0,1,\ldots,Q\}
\end{equation}
in clustered mode and
\begin{equation}
\mathcal{V} = \{0,1,\ldots,M\}
\end{equation}
in direct mode.

The clustering stage stores both a forward map and an inverse membership map. During runtime sampling, if a
clustered token $q > 0$ is drawn, one concrete dictionary row is chosen uniformly from the members of that
cluster. This preserves stochastic diversity while keeping the language-model vocabulary compact.

\subsection{LSTM Sequence Model}

Let $t_u \in \mathcal{V}$ denote the token at slot $u$. The knot generator models the autoregressive sequence
distribution
\begin{equation}
P(t_1,\ldots,t_T) = \prod_{u=1}^{T} P\!\left(t_u \mid t_1,\ldots,t_{u-1}\right).
\end{equation}
Each token is embedded into a continuous vector
\begin{equation}
e_u = E[t_u],
\end{equation}
where the embedding matrix $E$ is initialized from the normalized dictionary rows or cluster centroids. The
embedded sequence is processed by an LSTM,
\begin{equation}
(h_u, c_u) = \operatorname{LSTM}(e_u, h_{u-1}, c_{u-1}),
\end{equation}
followed by a linear prediction head
\begin{equation}
\ell_u = W h_u + b,
\qquad
P(t_{u+1} = q \mid t_{\le u}) = \operatorname{softmax}(\ell_u)_q.
\end{equation}

Training minimizes a weighted next-token cross-entropy:
\begin{equation}
\mathcal{L}
=
- \sum_{u=1}^{T-1}
w_{t_{u+1}}
\log P(t_{u+1} \mid t_{\le u}),
\end{equation}
where the no-knot token $0$ receives a reduced class weight. This is necessary because empty slots dominate the
longitudinal sequence distribution and would otherwise overwhelm the nonzero knot tokens.

At inference time, the generator samples a fresh sequence of length matching the board slot count. The sampler
supports both top-$k$ truncation and nucleus filtering. If a trained checkpoint is unavailable, the runtime can
fall back to a first-order Markov chain estimated from empirical token transitions in the training dataset.

\subsection{Runtime Decoding and Post-Processing}

After a token sequence has been sampled, token $i$ is assigned to the longitudinal slot
\begin{equation}
z_i = z_{\min} + i \Delta z,
\qquad i = 1,\ldots,N_{\mathrm{slot}}.
\end{equation}
Zero tokens are discarded. In clustered mode, every surviving token is decoded to a concrete dictionary row as
described above. Bad dictionary rows are removed, and the quartic diameter law is summarized by
\begin{equation}
L_{100} = a_1 100^4 + a_2 100^3 + a_3 100^2 + a_4 100.
\end{equation}
Only knots with
\begin{equation}
L_{100,\min} \le L_{100} \le L_{100,\max}
\end{equation}
are kept.

The simulator then derives the bump-law parameters through
\begin{equation}
B_{\mathrm{bump}} = 2,
\end{equation}
\begin{equation}
A_{\mathrm{bump}} = \left|0.0217\,L_{100} - 0.2\right|,
\end{equation}
\begin{equation}
A_{\mathrm{exp}} = 0.0056\,L_{100} + 1.96,
\end{equation}
and sets
\begin{equation}
k = 0.99, \qquad k_p = 0.95.
\end{equation}
For generated knots, the dead-zone width is repaired when necessary so that
\begin{equation}
R_D - R_L \ge \Delta_{\min}.
\end{equation}

Finally, the candidate knot set is checked for pairwise overlap by sampling each knot axis together with its
local diameter envelope. Two samples from different knots are deemed overlapping when
\begin{equation}
\left\lVert p_i - p_j \right\rVert_2 < r_i + r_j - \delta,
\end{equation}
where $p_i$ and $p_j$ are sampled axis points, $r_i$ and $r_j$ are the associated local radii, and $\delta$ is
the configured clearance. Sequences with overlapping knots are rejected and resampled until either a valid
candidate is found or the retry budget is exhausted.

\section{Growth-Layer Generator}

\subsection{Base Mesh and Nominal Growth Layers}

The simulator starts from a Cartesian grid
\begin{equation}
\mathcal{G} = \{(X_{ijk},Y_{ijk},Z_{ijk})\},
\end{equation}
constructed from one-dimensional coordinate arrays in the board domain. On this grid the transverse polar
coordinates are
\begin{equation}
\Theta = \operatorname{atan2}(Y,X),
\qquad
R = \sqrt{X^2 + Y^2}.
\end{equation}

For each growth-layer index $i$, the model stores a spline $s_i(\Theta)$ that defines the nominal layer radius:
\begin{equation}
R_{o,i}(\Theta) = s_i(\Theta).
\end{equation}
Without knots, crook, or taper, the implicit field of the $i$-th layer would be
\begin{equation}
g_i^{\mathrm{nom}}(X,Y,Z) = X^2 + Y^2 - R_{o,i}(\Theta)^2.
\end{equation}
The zero level set of this field is the nominal growth-layer surface.

\subsection{Crook and Taper}

The growth-layer generator does not depend on how crook and taper were sampled. It only assumes that the
current board realization provides two transverse displacement functions,
\[
c_x(Z), \qquad c_y(Z),
\]
and one taper function,
\[
t(Z).
\]
These functions may come from a random multi-component crook model, from manual parameters, or from legacy
quadratic terms. The growth-layer equations consume them only through
\begin{equation}
X_c = X + c_x(Z),
\qquad
Y_c = Y + c_y(Z),
\end{equation}
\begin{equation}
\widetilde{R}_{o,i}(\Theta,Z) = R_{o,i}(\Theta) - t(Z).
\end{equation}

Thus crook translates the transverse coordinates as a function of the longitudinal position, whereas taper
reduces the local effective ring radius.

\subsection{Rotation into the Knot Frame}

For a fixed knot $n$, the crook-corrected transverse point can be rotated directly in Cartesian form. The local
knot-frame coordinates are
\begin{equation}
\xi_n = X_c \cos\theta_{0,n} - Y_c \sin\theta_{0,n},
\qquad
\eta_n = X_c \sin\theta_{0,n} + Y_c \cos\theta_{0,n}.
\end{equation}
The coordinate $\xi_n$ is the radial coordinate in the rotated transverse plane, while $\eta_n$ is the
orthogonal tangential coordinate. This expression is exactly equivalent to first writing
$X_c=r\cos\phi$, $Y_c=r\sin\phi$ and then evaluating
\[
\xi_n = r\cos(\phi+\theta_{0,n}), \qquad
\eta_n = r\sin(\phi+\theta_{0,n}),
\]
but it avoids an unnecessary conversion to polar coordinates.

\subsection{Knot Axis and Influence Metric}

The knot centerline is modeled by a parabola in the $(\xi_n,Z)$ plane:
\begin{equation}
Z_{c,n}(\xi_n) = c_{1,n}\xi_n^2 + c_{2,n}\xi_n + z_{0,n}.
\end{equation}
The signed longitudinal offset from the knot axis is
\begin{equation}
\Delta Z_n = Z - Z_{c,n}(\xi_n).
\end{equation}

The anisotropic knot-deviation metric is then
\begin{equation}
p_n
=
\sqrt{
\Delta Z_n^2
+
\frac{1}{k_{p,n}^2}
\left(\xi_n^2 + \eta_n^2\right)
\operatorname{atan2}(\eta_n,\xi_n)^2
}.
\end{equation}
To avoid singular behavior near the knot axis, the implementation applies a smooth lower clamp:
\begin{equation}
\overline{p}_n
=
\sigma\!\left(\alpha (p_n - p_{\min})\right)p_n
+
\left[1 - \sigma\!\left(\alpha (p_n - p_{\min})\right)\right]p_{\min},
\end{equation}
where $\sigma(s) = (1 + e^{-s})^{-1}$ is the logistic sigmoid. The maximum active influence distance is
\begin{equation}
p_{\max,n}
=
\left(
\frac{
A_{\mathrm{bump},n}\,\widetilde{R}_{o,i}^{A_{\mathrm{exp},n}-1}
}{
1-k_n
}
\right)^{1/B_{\mathrm{bump},n}}.
\end{equation}

\subsection{Living-Knot Radius Law}

When the dead-knot extension is disabled, the knot affects the layer only in the forward half-space
$\xi_n \ge 0$ and only while $\overline{p}_n \le p_{\max,n}$. The effective layer radius for knot $n$ is
\begin{equation}
R_{i,n}
=
\begin{cases}
\widetilde{R}_{o,i}, &
\overline{p}_n > p_{\max,n} \ \text{or}\ \xi_n < 0, \\[4pt]
k_n \widetilde{R}_{o,i}
+
A_{\mathrm{bump},n}\widetilde{R}_{o,i}^{A_{\mathrm{exp},n}}
\overline{p}_n^{-B_{\mathrm{bump},n}}, &
\text{otherwise}.
\end{cases}
\end{equation}
If knot deviations are disabled entirely, then the simulator simply sets
\begin{equation}
R_{i,n} = \widetilde{R}_{o,i}.
\end{equation}

\subsection{Dead-Knot Transition}

When dead-knot logic is enabled, the model introduces a transition annulus
\begin{equation}
R_{L,n} \le \widetilde{R}_{o,i} \le R_{D,n}.
\end{equation}
The code defines
\begin{equation}
S_n = k_n R_{L,n},
\qquad
T_n = A_{\mathrm{bump},n} R_{L,n}^{A_{\mathrm{exp},n}},
\qquad
K_{\mathrm{red},n} = \frac{\widetilde{R}_{o,i} - R_{L,n}}{R_{D,n} - R_{L,n}},
\end{equation}
\begin{equation}
\Delta R_n = R_{D,n} - R_{L,n},
\qquad
\gamma = \frac{1}{1.5},
\end{equation}
\begin{equation}
\Delta p_n
=
\operatorname{Re}\!\left(
\sqrt{\Delta R_n^2 - \gamma^2\left(R_{D,n} - \widetilde{R}_{o,i}\right)^2}
\right)
-
\sqrt{1-\gamma^2}\,\Delta R_n,
\end{equation}
\begin{equation}
\overline{p}_{n,\mathrm{alt}} = \overline{p}_n + \Delta p_n.
\end{equation}
Within the transition region the effective radius becomes
\begin{equation}
R_{i,n}^{\mathrm{dead}}
=
K_{\mathrm{red},n} R_{D,n}
+
\left(1-K_{\mathrm{red},n}\right)
\left(
S_n + T_n \overline{p}_{n,\mathrm{alt}}^{-B_{\mathrm{bump},n}}
\right).
\end{equation}

The full dead-knot piecewise law is therefore
\begin{equation}
R_{i,n}
=
\begin{cases}
\widetilde{R}_{o,i}, &
\overline{p}_n > p_{\max,n}
\ \text{or}\ \xi_n < 0
\ \text{or}\ \widetilde{R}_{o,i} > R_{D,n}, \\[4pt]
R_{i,n}^{\mathrm{dead}}, &
R_{L,n} \le \widetilde{R}_{o,i} \le R_{D,n}, \\[4pt]
k_n \widetilde{R}_{o,i}
+
A_{\mathrm{bump},n}\widetilde{R}_{o,i}^{A_{\mathrm{exp},n}}
\overline{p}_n^{-B_{\mathrm{bump},n}}, &
\text{otherwise}.
\end{cases}
\end{equation}
Finally the implementation enforces
\begin{equation}
R_{i,n} \leftarrow \max(R_{i,n},1).
\end{equation}

\subsection{Growth-Layer Field and Combination of Knots}

For one layer $i$ and one knot $n$, the knot-specific implicit field is
\begin{equation}
g_{i,n}(X,Y,Z) = \xi_n^2 + \eta_n^2 - R_{i,n}^2.
\end{equation}
The actual growth-layer field is the lower envelope over all knots:
\begin{equation}
g_i(X,Y,Z) = \min_n g_{i,n}(X,Y,Z).
\end{equation}
The deformed growth-layer surface is the zero isosurface
\begin{equation}
\Gamma_i = \{(X,Y,Z) : g_i(X,Y,Z) = 0\}.
\end{equation}

If no knots are present, the simulator still produces a valid ring field:
\begin{equation}
g_i(X,Y,Z) = X_c^2 + Y_c^2 - \max(\widetilde{R}_{o,i},1)^2.
\end{equation}

When dead-knot mode is active, the implementation applies an additional nesting correction from the outside
inward:
\begin{equation}
\widetilde{g}_N = g_N,
\qquad
\widetilde{g}_i = \max(g_i, \widetilde{g}_{i+1}),
\qquad
i = N-1,\ldots,1.
\end{equation}
This prevents an inner layer from protruding outside the next outer layer in the dead-knot transition zone.

\subsection{Knot Scalar Field}

Independently of the ring field, the simulator constructs a knot-body field. The quartic diameter law
$L_n(\xi_n)$ is first evaluated. In dead-knot mode, the law is tapered after $R_{L,n}$ by freezing its value at
$R_{L,n}$ and multiplying it by a quadratic decay
\begin{equation}
h(t) = 1 - t^2,
\qquad
t = \frac{\xi_n - R_{L,n}}{R_{D,n} - R_{L,n}},
\end{equation}
with $L_n(\xi_n)=0$ for $\xi_n > R_{D,n}$.

The knot scalar field is
\begin{equation}
K_n(X,Y,Z)
=
\Delta Z_n^2 + \frac{\eta_n^2}{k_{p,n}^2}
- \left(\frac{L_n(\xi_n)}{2}\right)^2.
\end{equation}
The combined field is
\begin{equation}
K_{\min}(X,Y,Z) = \min_n K_n(X,Y,Z).
\end{equation}
By construction, $K_n < 0$ indicates the interior of the knot body, $K_n = 0$ is the nominal knot boundary,
and $K_n > 0$ is outside the knot.

For visualization and masking, the code also keeps separate live and dead subsets:
\begin{equation}
K_{n,\mathrm{live}} =
\begin{cases}
K_n, & \xi_n \le R_{L,n}, \\
\mathrm{NaN}, & \text{otherwise},
\end{cases}
\end{equation}
\begin{equation}
K_{n,\mathrm{dead}} =
\begin{cases}
K_n, & R_{L,n} < \xi_n \le R_{D,n}, \\
\mathrm{NaN}, & \text{otherwise}.
\end{cases}
\end{equation}

\subsection{Exported World-Space Knot Axes}

The implementation also exports explicit world-space knot axes. For a sampled radial coordinate $\xi$ with
$\eta = 0$, the axis centerline is
\begin{equation}
Z^{\mathrm{axis}}_n(\xi) = c_{1,n}\xi^2 + c_{2,n}\xi + z_{0,n},
\end{equation}
\begin{equation}
X^{\mathrm{axis}}_n(\xi) = \xi \cos\theta_{0,n},
\qquad
Y^{\mathrm{axis}}_n(\xi) = -\xi \sin\theta_{0,n}.
\end{equation}
Crook compensation then maps the axis into world coordinates:
\begin{equation}
X^{\mathrm{axis,world}}_n(\xi)
=
X^{\mathrm{axis}}_n(\xi) - c_x\!\left(Z^{\mathrm{axis}}_n(\xi)\right),
\end{equation}
\begin{equation}
Y^{\mathrm{axis,world}}_n(\xi)
=
Y^{\mathrm{axis}}_n(\xi) - c_y\!\left(Z^{\mathrm{axis}}_n(\xi)\right),
\end{equation}
\begin{equation}
Z^{\mathrm{axis,world}}_n(\xi) = Z^{\mathrm{axis}}_n(\xi).
\end{equation}
These are the quantities exported to the viewer and to downstream geometry consumers.

\section{Fiber Orientation Model}

\subsection{Starting Point: Foley's Fiber Paradigm}

The fiber model follows the same geometric idea as Foley's fiber paradigm. Foley first determines the
projection of the fiber direction in a longitudinal--tangential plane around a single knot. That projected
direction is then lifted into three dimensions by enforcing that the full fiber direction must lie in the
growth surface on which the material was produced.

The present simulator uses the same separation, but with the notation and coordinate convention used
throughout this manual. A board point is denoted by
\[
\bm{x}=(X,Y,Z),
\]
where $X$ is width, $Y$ is thickness, and $Z$ is board length. For any growth layer, the simulator already has
an implicit field $g_i(X,Y,Z)$ from the growth-layer model. The local radial material direction is therefore
the normalized growth-layer normal
\begin{equation}
\bm{n}
=
\frac{\nabla g_{i^\ast}}{\|\nabla g_{i^\ast}\|_2},
\qquad
i^\ast(X,Y,Z)=\operatorname*{argmin}_i |g_i(X,Y,Z)|.
\end{equation}
The fiber direction $\bm{f}$ must satisfy
\begin{equation}
\bm{f}\cdot\bm{n}=0.
\end{equation}
This is the direct analogue of Foley's orthogonality condition between the fiber direction and the normal to
the growth surface. The difference is that Foley describes the growth surface with an explicit radial function,
whereas the board generator uses the already-computed implicit fields $g_i$.

In Foley's single-knot derivation, the knot plane is aligned with the global axes. Therefore the projected
fiber pattern can be written directly in the thesis coordinates. In the board generator, each knot has its own
azimuth, the pith may be crooked, and many knots may be present in the same board. The projected fiber pattern
must therefore be evaluated in each knot's local frame and then mapped back to the board frame.

\subsection{Board Coordinates and Knot Coordinates}

The same crook-corrected coordinates used by the growth-layer model are used before evaluating the fiber
disturbance:
\begin{equation}
X_c = X + c_x(Z),
\qquad
Y_c = Y + c_y(Z).
\end{equation}
For knot $n$, with azimuth $\theta_{0,n}$, the transverse point is rotated into the knot frame as
\begin{equation}
\xi_n = X_c \cos\theta_{0,n} - Y_c \sin\theta_{0,n},
\qquad
\eta_n = X_c \sin\theta_{0,n} + Y_c \cos\theta_{0,n}.
\end{equation}
Here $\xi_n$ is the coordinate along the radial knot plane and $\eta_n$ is the tangential coordinate around the
stem. The knot axis in this frame is
\begin{equation}
Z_{c,n}(\xi_n)=c_{1,n}\xi_n^2+c_{2,n}\xi_n+z_{0,n}.
\end{equation}
It is useful in the fiber section to name the signed longitudinal distance from the point to the knot axis:
\begin{equation}
\zeta_n = Z - Z_{c,n}(\xi_n).
\end{equation}
This $\zeta_n$ is the same quantity denoted $\Delta Z_n$ in the growth-layer section. In Foley's notation, the
corresponding coordinate is the longitudinal coordinate measured relative to the knot centerline.

The associated knot-frame basis vectors, expressed in board coordinates, are
\begin{equation}
\bm{e}_{\xi,n}=(\cos\theta_{0,n},-\sin\theta_{0,n},0),
\qquad
\bm{e}_{\eta,n}=(\sin\theta_{0,n},\cos\theta_{0,n},0),
\qquad
\bm{e}_Z=(0,0,1).
\end{equation}
Foley did not need this rotation step because the thesis derivation considered a single knot whose knot plane
was already aligned with the chosen coordinate system.

\subsection{Projected Flow Pattern Around One Knot}

For a fixed value of $\xi_n$, the simulator considers the plane spanned by $\eta_n$ and the local longitudinal
offset $\zeta_n$. This is the board-generator counterpart of Foley's longitudinal--tangential plane. The knot
diameter in that plane is
\begin{equation}
L_n(\xi_n)
=
a_{1,n}\xi_n^4+a_{2,n}\xi_n^3+a_{3,n}\xi_n^2+a_{4,n}\xi_n.
\end{equation}
The projected fiber path is represented by a Rankine-oval stream pattern around an effective obstacle of
width
\begin{equation}
b_n(\xi_n)=k_{p,n}L_n(\xi_n).
\end{equation}
Thus $L_n$ controls the longitudinal size of the obstacle and $b_n$ controls the tangential size. This matches
the role of Foley's longitudinal and transverse knot diameters, but the values are evaluated from the
runtime knot law at the local coordinate $\xi_n$.

The flow construction is parameterized by a source--sink half-spacing $a_n$ and a strength parameter $G_n$.
The implementation initializes $a_n$ by
\begin{equation}
a_n^{\mathrm{approx}}
=
\frac{\sqrt{\nu}\,L_n(\xi_n)\,[b_n(\xi_n)/2]}
{\sqrt{(2\nu+1)L_n(\xi_n)^2 + (4\nu-4)[b_n(\xi_n)/2]^2}},
\qquad
\nu=7.
\end{equation}
When the exact mode is selected, this initial value is refined by Newton iterations for the residual
\begin{equation}
F(a_n)
=
8a_n y_n
-
\left[
\operatorname{atan2}(y_n,-a_n)-\operatorname{atan2}(y_n,a_n)
\right]
\left(L_n(\xi_n)^2-4a_n^2\right),
\qquad
y_n=\frac{b_n(\xi_n)}{2}.
\end{equation}
After $a_n$ is known, the strength parameter is
\begin{equation}
G_n
=
\frac{\pi L_n(\xi_n)^2 - 4\pi a_n^2}{4a_n}.
\end{equation}

Let $\psi_n(\eta_n,\zeta_n)$ be the local stream function. The implementation evaluates the derivatives
\begin{equation}
\frac{\partial \psi_n}{\partial \eta_n}
=
\frac{G_n}{2\pi}
\left[
\frac{a_n + \zeta_n}{\eta_n^2 + (a_n + \zeta_n)^2}
+
\frac{a_n - \zeta_n}{\eta_n^2 + (a_n - \zeta_n)^2}
\right]
+1,
\end{equation}
\begin{equation}
\frac{\partial \psi_n}{\partial \zeta_n}
=
-\frac{G_n}{2\pi}
\left[
\frac{\eta_n}{\eta_n^2 + (a_n + \zeta_n)^2}
-
\frac{\eta_n}{\eta_n^2 + (a_n - \zeta_n)^2}
\right].
\end{equation}
The projected fiber direction is tangent to the streamlines. Therefore the unnormalized projected components
are
\begin{equation}
u_{\eta,n}=-\frac{\partial \psi_n}{\partial \zeta_n},
\qquad
u_{Z,n}=\frac{\partial \psi_n}{\partial \eta_n}.
\end{equation}
The normalized projected direction is
\begin{equation}
\widehat{u}_{\eta,n}
=
\frac{u_{\eta,n}}{\sqrt{u_{\eta,n}^2+u_{Z,n}^2}},
\qquad
\widehat{u}_{Z,n}
=
\frac{u_{Z,n}}{\sqrt{u_{\eta,n}^2+u_{Z,n}^2}}.
\end{equation}
This gives the same kind of projected direction that Foley denotes by the projection of the fiber vector in the
longitudinal--tangential plane. The board generator does not model separate $\alpha$- and $\beta$-patterns in
the current implementation; it uses this single projected flow field for the clear-wood fiber deviation around
the knot.

\subsection{Lifting the Projected Direction into Three Dimensions}

The projected direction is first mapped from the knot frame to the board frame:
\begin{equation}
\bm{t}^{(0)}_n
=
\widehat{u}_{\eta,n}\bm{e}_{\eta,n}
+
\widehat{u}_{Z,n}\bm{e}_Z.
\end{equation}
Componentwise, this is
\begin{equation}
t^{(0)}_{x,n}=\widehat{u}_{\eta,n}\sin\theta_{0,n},
\qquad
t^{(0)}_{y,n}=\widehat{u}_{\eta,n}\cos\theta_{0,n},
\qquad
t^{(0)}_{z,n}=\widehat{u}_{Z,n}.
\end{equation}

At this point $\bm{t}^{(0)}_n$ is only the projected flow direction. To obtain a three-dimensional fiber
direction, the solver adds the missing radial component in the local $\xi_n$ direction:
\begin{equation}
\bm{t}_n
=
\bm{t}^{(0)}_n+\delta_n\bm{e}_{\xi,n}.
\end{equation}
The scalar $\delta_n$ is chosen so that the final direction is tangent to the nearest growth layer:
\begin{equation}
\bm{n}\cdot
\left(
\bm{t}^{(0)}_n+\delta_n\bm{e}_{\xi,n}
\right)
=0.
\end{equation}
Solving this one linear equation gives
\begin{equation}
\delta_n
=
-\frac{
n_x t^{(0)}_{x,n}
+n_y t^{(0)}_{y,n}
+n_z t^{(0)}_{z,n}
}{
n_x\cos\theta_{0,n}-n_y\sin\theta_{0,n}+\varepsilon
},
\end{equation}
where $\varepsilon$ is a small numerical stabilizer. The vector $\bm{t}_n$ is then normalized. This is the
board-generator version of Foley's step in which the unknown radial component is recovered from the condition
that the full fiber vector be orthogonal to the growth-surface normal.

\subsection{Multiple Knots and Default Clear Wood}

Foley's thesis develops the fiber paradigm for one knot. The board generator evaluates the same construction
for every knot $n$ that exists in the board. If more than one knot gives a valid candidate direction at the
same point, the implementation selects the candidate with the largest absolute longitudinal component:
\begin{equation}
n^\ast
=
\operatorname*{argmax}_n |t_{z,n}|.
\end{equation}
The clear-wood fiber direction is then
\begin{equation}
\bm{f}_{\mathrm{clear}}(X,Y,Z)
=
\bm{t}_{n^\ast}(X,Y,Z).
\end{equation}
If the projected-flow calculation is invalid or undefined at a point, the simulator falls back to the straight
longitudinal direction
\begin{equation}
\bm{f}_{\mathrm{default}}=(0,0,1).
\end{equation}

\subsection{Optional In-Knot Axis Override}

The projected-flow construction describes how stem fibers deviate around the knot. Inside the knot body itself,
the simulator can optionally replace this clear-wood direction by the tangent to the knot axis. The knot-body
region is defined by the scalar field $K_n$ from the growth-layer section; the override is active where
\begin{equation}
K_n(X,Y,Z)<K_{\mathrm{in}},
\end{equation}
with $K_{\mathrm{in}}$ set by the configured interior threshold.

Because the local knot axis is
\[
Z_{c,n}(\xi_n)=c_{1,n}\xi_n^2+c_{2,n}\xi_n+z_{0,n},
\]
its tangent in the knot frame is
\begin{equation}
\bm{\tau}^{\mathrm{kf}}_n
=
\frac{(1,0,2c_{1,n}\xi_n+c_{2,n})}
{\sqrt{1+(2c_{1,n}\xi_n+c_{2,n})^2}}.
\end{equation}
Rotating this tangent into board coordinates gives
\begin{equation}
\bm{\tau}_n
=
\left(
\frac{\cos\theta_{0,n}}{\sqrt{1+(2c_{1,n}\xi_n+c_{2,n})^2}},
-\frac{\sin\theta_{0,n}}{\sqrt{1+(2c_{1,n}\xi_n+c_{2,n})^2}},
\frac{2c_{1,n}\xi_n+c_{2,n}}{\sqrt{1+(2c_{1,n}\xi_n+c_{2,n})^2}}
\right).
\end{equation}
This override is separate from Foley's projected-flow calculation. Two implementation options modify where or
how it is applied:
\begin{enumerate}
\item The override can be suppressed in the dead-knot region by excluding points where $K_{n,\mathrm{dead}}$ is
finite.
\item The override direction can be sign-reversed for $\zeta_n>0$ so that the axis tangent flips on the
positive longitudinal side of the knot.
\end{enumerate}

The resulting field is
\begin{equation}
\bm{f}(X,Y,Z)
=
\begin{cases}
\bm{\tau}_n, & \text{inside the active knot-override region},\\[4pt]
\bm{f}_{\mathrm{clear}}, & \text{otherwise, when a valid knot-induced direction exists},\\[4pt]
(0,0,1), & \text{otherwise}.
\end{cases}
\end{equation}

\subsection{Summary of Differences from Foley's Single-Knot Model}

The board-generator fiber model should be read as a direct extension of Foley's construction rather than a
different theory. The projected fiber direction is still obtained from a Rankine-oval flow pattern in a
longitudinal--tangential plane, and the missing radial component is still determined by the requirement that
the fiber lie in the local growth surface.

The added steps are implementation necessities. The point is first transformed into the crook-adjusted
transverse coordinate system and rotated into each knot frame because the simulated board can contain knots at
arbitrary azimuths. The projected direction is then rotated back to the board frame. Finally, the simulator
resolves competition among multiple knots and can optionally replace the clear-wood direction inside the knot
by the knot-axis tangent. These additions are the parts absent from Foley's aligned, single-knot coordinate
system.

\end{document}
